\section{总结与延伸}

\subsection{讲者的核心总结}

\begin{figure}[H]
\centering
\includegraphics[width=0.95\textwidth]{slides-images/slide-064.jpg}
\caption{全课总结 Slides\protect\footnotemark}
\end{figure}
\footnotetext{Slides 第65页（最后一页）。}

Yann Dubois 在课程结尾给出了一条最重要的建议：

\begin{importantbox}{最好的评估方法：看你的输出}
不要盲信数字。太多人只看基准分数就下结论。在 Alpaca 项目中，研究团队一开始也只看 AlpacaEval 分数，但真正让他们意识到模型质量的是\textbf{实际使用和查看模型输出}。一个在学术基准上分数很低的模型，实际使用中可能表现不错；反之亦然。

\textbf{``The best evaluation is just check your outputs.''}
\end{importantbox}

\subsection{全课知识图谱}

\begin{center}
\begin{tikzpicture}[
    node distance=0.6cm,
    mainnode/.style={draw, rounded corners, minimum width=4cm, minimum height=0.7cm, font=\small, fill=blue!5},
    subnode/.style={draw=none, text width=6cm, font=\scriptsize, align=left},
    arr/.style={-{Stealth[length=3mm]}, thick}
]
\node[mainnode] (why) {为什么需要评估};
\node[mainnode, below=of why] (close) {封闭式评估};
\node[mainnode, below=of close] (open) {开放式评估};
\node[mainnode, below=of open] (llm) {LLM 作为评估者};
\node[mainnode, below=of llm] (curr) {当前 LLM 评估实践};
\node[mainnode, below=of curr] (issue) {问题与挑战};

\draw[arr] (why) -- (close);
\draw[arr] (close) -- (open);
\draw[arr] (open) -- (llm);
\draw[arr] (llm) -- (curr);
\draw[arr] (curr) -- (issue);

\node[subnode, right=0.8cm of why] {训练/开发/模型选择/部署/发表};
\node[subnode, right=0.8cm of close] {Accuracy/F1、SuperGLUE、虚假相关性};
\node[subnode, right=0.8cm of open] {BLEU/ROUGE、BERTScore、BLEURT};
\node[subnode, right=0.8cm of llm] {AlpacaEval、Chatbot Arena、偏差-方差权衡};
\node[subnode, right=0.8cm of curr] {困惑度、MMLU/HELM、竞技场式};
\node[subnode, right=0.8cm of issue] {一致性、污染、单一化、偏见放大};
\end{tikzpicture}
\end{center}

\subsection{关键 Takeaways}

\begin{importantbox}{八条核心原则}
\begin{enumerate}
    \item \textbf{不同阶段需要不同的评估}：训练需要快速可微指标，部署需要可信的绝对指标
    \item \textbf{封闭式评估并不简单}：指标选择、标签质量、虚假相关性都需要仔细考虑
    \item \textbf{开放式评估是 NLP 特有的难题}：没有唯一正确答案，答案质量是连续谱
    \item \textbf{BLEU/ROUGE 已经过时}：无法捕捉语义相似性，但因学术惯性仍在广泛使用
    \item \textbf{人工评估是金标准但代价高昂}：标注者不一致、不可复现、激励不对齐
    \item \textbf{LLM 评估是革命性突破}：快 100 倍、便宜 100 倍，一致性甚至超过人类
    \item \textbf{数据污染是隐患}：特别是对闭源模型，私有测试集和动态基准是应对策略
    \item \textbf{永远要看实际输出}：不要盲信任何基准分数——实际使用模型是最好的评估
\end{enumerate}
\end{importantbox}

\subsection{拓展阅读}

\begin{itemize}
    \item Papineni et al., BLEU: a Method for Automatic Evaluation of Machine Translation, ACL 2002
    \item Lin, ROUGE: A Package for Automatic Evaluation of Summaries, 2004
    \item Zhang et al., BERTScore: Evaluating Text Generation with BERT, ICLR 2020
    \item Sellam et al., BLEURT: Learning Robust Metrics for Text Generation, ACL 2020
    \item Dubois et al., AlpacaFarm: A Simulation Framework for Methods that Learn from Human Feedback, NeurIPS 2023
    \item Zheng et al., Judging LLM-as-a-Judge with MT-Bench and Chatbot Arena, NeurIPS 2023
    \item Hendrycks et al., Measuring Massive Multitask Language Understanding, ICLR 2021
    \item Gururangan et al., Annotation Artifacts in Natural Language Inference Data, NAACL 2018
    \item Chris Potts's CS224U 课程：\url{https://web.stanford.edu/class/cs224u/}（评估指标详细讲解）
    \item Chatbot Arena 排行榜：\url{https://chat.lmsys.org/}
    \item AlpacaEval 排行榜：\url{https://tatsu-lab.github.io/alpaca_eval/}
\end{itemize}

\end{document}
